\documentclass[12pt]{article}

%Lei de Ampere para um condutor cilindrico infinito
%Faro, Tue Mar 11 22:54:20 WET 2014
%Written by Tordar

\usepackage{graphicx}
\usepackage{pst-all}



\begin{document}

\pagestyle{empty}

\begin{pspicture}(-6,-4)(6,4)
\pscircle[linewidth=1.5mm,linecolor=black,fillstyle=vlines](-4,0){3}
\pscircle[linewidth=1.5mm,linecolor=black,fillstyle=vlines]( 4,0){3}
\pscircle[linewidth=0.5mm,linecolor=black,linestyle=dashed,dash=4mm 2mm](-4,0){2}
\pscircle[linewidth=0.5mm,linecolor=black,linestyle=dashed,dash=4mm 2mm]( 4,0){4}
\psline[linewidth=1mm,linecolor=blue,arrows=->](-4,0)(-1,0)
\psline[linewidth=1mm,linecolor=blue,arrows=->]( 4,0)( 7,0)
\psline[linewidth=1mm,linecolor=blue,arrows=->](-4,0)(-4,2)
\psline[linewidth=1mm,linecolor=blue,arrows=->]( 4,0)( 4,4)
\rput{0}( -3,-0.5){\scalebox{1.5}{\psframebox*{$R$}}}
\rput{0}(5.5,-0.5){\scalebox{1.5}{\psframebox*{$R$}}}
\rput{0}(-4.5,1){\scalebox{1.5}{\psframebox*{$r$}}}
\rput{0}( 3.5,1){\scalebox{1.5}{\psframebox*{$r$}}}
\rput{0}(-4,2.4){\scalebox{1.5}{\psframebox*{$L$}}}
\rput{0}( 4,4.4){\scalebox{1.5}{$L$}}
\end{pspicture}

\end{document}
